%%%%%%%%%%%%%%%%%%%%%%%%%%%%%%%%%%%% Chapter Template

 \chapter{Introduction} 	% Main chapter title
 \label{Chapter1} 		% For referencing the chapter elsewhere, usage \ref{Chapter1}

% %%%%%%%%%%%%%%%%%%%%%%%%%%%%%%%%%%%%

% This document aims to guide the Student in the preparation of the Project/Internship report, for the 3rd year of the \ac{leec} and \ac{leti}, from the \ac{dee}, of the \ac{isep}.

% O autor deverá ter em consideração as seguintes regras gerais na elaboração do documento:

% \begin{itemize}
% 	\item É fundamental refletir, antes de escrever, sobre a substância do que se pretende transmitir.
% 	\item Deve organizar o texto evitando a excessiva divisão do mesmo em tópicos que supostamente se enquadram no tema principal da secção a que pertencem. A excessiva especialização pode ser reveladora de falta de conhecimento e/ou reflexão. Neste sentido, deve, antes de iniciar a escrita, exercitar-se em refinar a organização do texto com o intuito de evitar que existam mais do que 2 níveis de ``profundidade'' em cada secção (respetivamente, subsecção e sub-subsecção). Note que o número de capítulos, secções e subsecções do presente documento não é vinculativo nem mesmo indicativo, apenas serve os propósitos do mesmo.
% 	\item O documento deve ser redigido em português ou inglês com um estilo adequado (evite o tom coloquial, lugares comuns e chavões) e correto do ponto de vista gramatical (quer do ponto de vista sintático quer semântico).
% 	\item Tenha especial cuidado com o uso de adjetivos (facilmente conduzem ao exagero), advérbios (nada, ou quase nada, acrescentam) e sinais de pontuação (em especial o uso correto das vírgulas).
% 	\item O estilo adotado para a redação deve ser coerente com as exigências de um trabalho científico encontrado em publicações impressas.
% 	\item De uma forma genérica deve usar a 3ª pessoa do singular (eventualmente do plural), exceção feita aos locais onde tal é claramente desajustado, por exemplo, na secção dos agradecimentos.
% 	\item Usar o estilo \textit{itálico} sempre que são utilizados termos em línguas diferentes da língua adotada no relatório.
% 	\item O uso de acrónimos implica que na 1ª vez que são utilizados se apresentem por extenso, colocando entre parênteses a respetiva sigla que se passará a usar. No entanto, é sempre possível, mais à frente no texto, e por uma questão de legibilidade, repetir o significado do acrónimo. Todos os acrónimos devem ser apresentados por ordem alfabética na secção ``Lista de Acrónimos''.
% 	\item O uso correto de unidades, seus múltiplos e submúltiplos.
% 	\item As imagens e tabelas devem, por princípio, aparecer no topo ou no fundo da página. A legendas surgem imediatamente após as figuras e listagens. No caso das tabelas as legendas antecedem as mesmas.
% 	\item Todas as figuras, tabelas e restantes listagens devem ser mencionadas no texto por forma a que fiquem enquadradas nas ideias transmitidas pelo autor. Esta referência, regra geral, deverá ser feita antes da ocorrência da figura, tabela ou listagem.
% 	\item Deve indicar ao longo do texto as referências documentais usadas, em especial nas citações (puras ou literais), assinaladas com a utilização de aspas, como também no caso de reutilização de gráficos, figuras, tabelas, fórmulas, etc., de outras fontes.
% \end{itemize}

% De uma forma já mais específica, neste primeiro capítulo obrigatório (``Introdução'') o autor deve:

% \begin{itemize}
% 	\item contextualizar a proposta de trabalho no âmbito da empresa, de um outro trabalho já realizado, do ponto de vista científico e/ou tecnológico, etc.,
% 	\item apresentar de forma clara os objetivos que se propõe atingir,
% 	\item descrever de forma sucinta, mas objetiva, a solução preconizada ou a hipótese colocada,
% 	\item apresentar de forma resumida, mas clara, os desenvolvimentos efetuados,
% 	\item identificar como foi validada e avaliada a solução encontrada,
% 	\item descrever a organização do documento.
% \end{itemize}

 	

%%%%%%%%%%%%%%%%%%%%%%%%%%%%%%%%%%%% SECTION 1

\section{State of the Art/Contextualization}	%Section title
\label{sec:Ch1.1}

Stethoscopes remain the most fundamental and prevasive diagnostic tools in clinical practice for the auscultation of cardiopulmonary and vascular sounds. Currently, traditional acoustic stethoscopes dominate over 90\% of the global market due to their high durability, low cost, and independence from electrical power sources \cite{market_report_2022}. However, purely acoustic auscultation is inherently limited not only by the fallibility of human auditory perception but also by the unrecognized physical degradation of the instrument over time. Studies have demonstrated that general practitioners miss up to 57\% of significant valvular heart diseases when relying on standard stethoscopes \cite{gardezi2017cardiac}. 

While frequently attributed to human auditory limitations, this high error rate may be significantly composed by unnoticed material failures; such as compromised acoustic seals in the eartips, hardened tubing, or worn diaphragms. These physical degradations silently alter the stethoscope's frequency response, attenuating vital low-frequency biological sounds before they ever reach the clinician's ear. To mitigate these human diagnostic errors, the medical field is experiencing the gradual emergence of electronic stethoscopes designed to amplify, filter, and digitize biological sounds.

This technological shift, while clinically beneficial, introduces a critical new engineering problem: these electronic devices actively alter and process the raw audio signal. Because both the acoustic	 integrity of traditional tubing and the digital processing of modern devices directly impact diagnostic accuracy, the development of a precise calibration system to verify signal fidelity has never been more urgent.

%%%%%%%%%%%%%%%%%%%%%%%%%%%%%%%%%%%% SECTION 2

\section{Project Description}
\label{sec:Ch1.2}

This project proposes the development of an acoustic calibration system designed to characterize the frequency response, and signal fidelity of stethoscopes. The proposed architecture consists of two primary domains: a physical hardware module for sound generation and detection, and a software platform for control, data acquisition, and spectral analysis.

%%%%%%%%%%%%%%%%%%%%%%%%%%%%%%%%%%%% SUBSECTION 1

\subsection{Objectives}
\label{sub:Ch1.2.1}

By the end of the project, the developed system should be able to:
\begin{itemize}
	\item \textbf{Simulate Cardiopulmonary Acoustics:} Generate controlled acoustic signals that mimic the frequency spectrum of biological sounds, particularly in the frequency range for cardiopulmonary auscultation.
	\item \textbf{Capture High-Fidelity Audio Data:} Accurately measure the stethoscope's output signal using a high-fidelity microphone and data acquisition system, ensuring that the measurement setup itself does not introduce significant distortions.
	\item \textbf{Analyze Frequency Response:} Analyze the acquired audio data to compute the stethoscope's frequency response curve utilizing python-based spectral analysis techniques.
	\item \textbf{Establish a Repeatable Protocol:} Develop a standardized testing protocol that can be easily replicated for routine stethoscope calibration in clinical settings, ensuring consistent and reliable performance assessments.
\end{itemize}



% you could use the pgfgantt package: https://ctan.org/pkg/pgfgantt



